% Generated draft: requires mathematical and source review.
\documentclass[11pt]{article}
\usepackage[margin=25mm]{geometry}
\usepackage{fontspec}
\setmainfont{Latin Modern Roman}
\usepackage{amsmath,amssymb,mathtools,unicode-math}
\setmathfont{Latin Modern Math}
\usepackage{enumitem,needspace,xurl,hyperref}
\setlength{\parindent}{0pt}
\setlength{\parskip}{5pt}
\setlength{\emergencystretch}{3em}
\newcommand{\N}{\mathbb{N}}
\newcommand{\Z}{\mathbb{Z}}
\newcommand{\Q}{\mathbb{Q}}
\newcommand{\R}{\mathbb{R}}
\newcommand{\C}{\mathbb{C}}
\newcommand{\F}{\mathbb{F}}
\newcommand{\A}{\mathbb{A}}
\newcommand{\PP}{\mathbb{P}}
\newcommand{\E}{\mathbb{E}}
\newcommand{\eps}{\varepsilon}
\newcommand{\dd}{\,\mathrm{d}}
\begin{document}
\section*{UnsolvedMath ID: 6800009}
\textbf{Source problem code:} catalogue entry
\subsection{Definitions and mathematical statement}
Let $G$ be a compact connected Lie group and $g$ a left-invariant Riemannian metric, so $L_a^*g=g$ for every $a\in G$, where $L_a(x)=ax$. The metric is bi-invariant if also $R_a^*g=g$ for every $a$, where $R_a(x)=xa$.

For a geodesic segment $\gamma:[0,T]\to G$, interpret its index as the fixed-endpoint Morse index of the energy: the maximal dimension of a subspace on which the index form
\[
 I_\gamma(V,V)=\int_0^T\bigl(
 |D_tV|^2-\langle R(V,\dot\gamma)\dot\gamma,V\rangle
 \bigr)\,\mathrm dt
\]
is negative definite, for $H^1$ vector fields along $\gamma$ vanishing at both endpoints, with the usual compatible curvature sign convention. Zero directions, including endpoint degeneracy, contribute nullity rather than index.

Gorodski's question asks whether
\[
 \left(\forall\gamma\text{ geodesic segment},\quad
        \operatorname{ind}(\gamma)\in2\Z\right)
 \quad\Longrightarrow\quad
 g\text{ is bi-invariant}.
\]
The parity condition concerns every segment and every segment length, not only closed geodesics or one chosen endpoint. No simple-connectivity or semisimplicity hypothesis is added. The fixed-endpoint interpretation is standard but is not explicitly defined in the short source statement.
\subsection{Short English statement}
Does even Morse index for every geodesic segment characterize bi-invariant metrics among left-invariant metrics on compact connected Lie groups? The source asks for the implication to bi-invariance, without additional group restrictions.

The supplied summary reports no verified direct resolution and cites related rigidity and spectral-isolation results as context only. It explicitly says those do not settle the even-index characterization. The OPEN-TRIAGE label and 2018 question are retained with current status marked for review.
\subsection{Sources}
\begin{itemize}
\item Local UnsolvedMath registry, permanent ID 6800009, source code catalogue entry, ``Bi-invariant metrics and multiplicity of conjugate points,'' proposed by Claudio Gorodski, University of São Paulo, Brazil, 2018. Dataset curation and metadata: Ulam AI and UnsolvedMath Contributors, CC BY 4.0.
\item Morgan and Pansu, \emph{A List of Open Problems in Differential Geometry} (2018), Question environment 9, \url{https://www.imo.universite-paris-saclay.fr/~pansu/problems_MTDG.tex}; source-TeX extraction, accessed 2026-07-29.
\item Supplied research-summary context: a SIGMA 2020 note on Gromov rigidity of bi-invariant metrics; ``Spectral isolation of bi-invariant metrics on compact Lie groups.'' No complete authors, publication details, or exact theorem statements are supplied for these references.
\item Source-specific provenance: the item is labelled a question, current status NEEDS\_REVIEW; the public source TeX is on a coauthor's publication page, with redistribution terms NEEDS\_REVIEW. These qualifications are separate from the dataset licence.
\end{itemize}
\end{document}
